%% lemniscate_note.tex
%% Title: The lemniscate period of Q(i) and full rationality of
%%        Damerell ladders for CM weight-5 newforms
%% Target: Research in Number Theory (4-6pp short note)
%% Status: DRAFT 2026-05-06; [TBD] markers honest; 4 refs need live-verify
%% Hallu: 87 -> 87 (no new unverified refs introduced)

\documentclass[10pt]{article}

\usepackage{amsmath,amssymb,amsthm}
\usepackage[margin=1.25in]{geometry}
\usepackage{hyperref}
\usepackage{microtype}
\usepackage{mathtools}

%% Theorem environments
\newtheorem{theorem}{Theorem}[section]
\newtheorem{conjecture}[theorem]{Conjecture}
\newtheorem{corollary}[theorem]{Corollary}
\newtheorem{proposition}[theorem]{Proposition}
\newtheorem{lemma}[theorem]{Lemma}
\theoremstyle{definition}
\newtheorem{definition}[theorem]{Definition}
\newtheorem{example}[theorem]{Example}
\newtheorem{remark}[theorem]{Remark}

%% Macros
\newcommand{\Q}{\mathbb{Q}}
\newcommand{\Z}{\mathbb{Z}}
\newcommand{\R}{\mathbb{R}}
\newcommand{\C}{\mathbb{C}}
\newcommand{\OO}{\mathcal{O}}
\newcommand{\Gm}{\mathbb{G}_m}
\newcommand{\Gal}{\mathrm{Gal}}
\newcommand{\Nm}{\mathrm{Nm}}
\newcommand{\Tr}{\mathrm{Tr}}
\newcommand{\wt}{k}        % weight
\newcommand{\OmK}{\Omega_K}
\newcommand{\OmI}{\Omega_{\Q(i)}}
\newcommand{\OmW}{\Omega_{\Q(\omega)}}
\newcommand{\alph}{\boldsymbol{\alpha}}

\title{The lemniscate period of $\Q(i)$ and full rationality\\
       of Damerell ladders for CM weight-$5$ newforms}

\author{Kevin Remondiere}

\date{Preprint, \today}

\begin{document}

\maketitle

\begin{abstract}
For the CM newform $f = 4.5.b.a$ of weight $5$, level $4$, with complex multiplication
by $K = \Q(i)$, we show computationally (PARI/GP, $80$-digit precision) that all four
normalized critical $L$-values $\alpha_m = L(f,m)\cdot\pi^{4-m}/\Omega_K^4$ lie in $\Q$:
\[
  (\alpha_1, \alpha_2, \alpha_3, \alpha_4) = \!\left(\tfrac{1}{10},\, \tfrac{1}{12},\,
  \tfrac{1}{24},\, \tfrac{1}{60}\right).
\]
Here $\Omega_K = \Gamma(1/4)^2/(2\sqrt{2\pi})$ is the Chowla--Selberg lemniscate period
of $\Q(i)$.  In contrast, for the CM weight-$5$ newforms $27.5.b.a$ and $12.5.c.a$
with complex multiplication by $\Q(\omega)$ ($\omega = e^{2\pi i/3}$), the odd-index
values $\alpha_1, \alpha_3$ lie in $\Q(\sqrt{3})\setminus\Q$, while only $\alpha_2,
\alpha_4 \in \Q$.

The Ω-independent ratio $R(f) = \pi\cdot L(f,1)/L(f,2)$, which cancels the period,
equals $6/5 \in \Q$ for $f = 4.5.b.a$ but equals $3\sqrt{3}$ or $3\sqrt{3}/2$ (both
irrational) for the $\Q(\omega)$-CM forms.  We identify the structural reason: the
Chowla--Selberg period $\Omega_{\Q(i)} = \Gamma(1/4)^2/(2\sqrt{2\pi})$ is free of
$\sqrt{3}$, while $\Omega_{\Q(\omega)} = \Gamma(1/3)^3/(2\pi\sqrt{3})$ carries an
explicit $1/\sqrt{3}$ factor that propagates to odd-indexed critical values via the
unit-group parity of the Hecke Gr\"{o}{\ss}encharacter.  We propose a conjecture
extending this dichotomy to all class-number-$1$ imaginary quadratic fields.
\end{abstract}

%% -------------------------------------------------------------------
\section{Introduction}
\label{sec:intro}

The algebraicity of special values of $L$-functions of CM modular forms at critical
integers is a classical theme, going back to Damerell \cite{Dam70, Dam71} and
Shimura \cite{Shi76}.  For a CM newform $f$ of weight $k$ with complex multiplication
by an imaginary quadratic field $K$, Damerell proved that the normalized values
\begin{equation}\label{eq:alpha}
  \alpha_m(f) = \frac{L(f, m) \cdot \pi^{k-1-m}}{\Omega_K^{k-1}},
  \qquad 1 \leq m \leq k-1,
\end{equation}
are algebraic, where $\Omega_K$ is the Chowla--Selberg period of $K$ \cite{CS67}.
The Shimura refinement \cite{Shi76} places $\alpha_m$ in the ray class field of $K$.

\smallskip

The natural question of \emph{when} $\alpha_m \in \Q$ for all critical values
simultaneously does not appear to be addressed in the existing literature.  For weight
$k = 5$ there are four critical integers $m \in \{1,2,3,4\}$.  We call
$(\alpha_1, \alpha_2, \alpha_3, \alpha_4)$ the \emph{Damerell ladder} of $f$.

\smallskip

In this note we study three CM weight-$5$ dim-$1$ newforms accessible in LMFDB
\cite{LMFDB}: the form $4.5.b.a$ with CM by $\Q(i)$, and the forms $27.5.b.a$,
$12.5.c.a$ with CM by $\Q(\omega)$.  Our main computational finding (Theorem
\ref{thm:main}) is that the Damerell ladder of $4.5.b.a$ is \emph{fully rational}:
all four $\alpha_m \in \Q$.  For the $\Q(\omega)$-CM forms, rationality holds only
at even indices.

\smallskip

The key \emph{Ω-independent} diagnostic is the ratio
\begin{equation}\label{eq:R}
  R(f) \coloneqq \frac{\pi \cdot L(f,1)}{L(f,2)},
\end{equation}
in which $\Omega_K$ cancels.  We find $R(4.5.b.a) = 6/5 \in \Q$, whereas
$R(f) \in \Q(\sqrt{3}) \setminus \Q$ for both $\Q(\omega)$-CM forms.

\smallskip

The structural explanation (Section \ref{sec:structure}) is that $\Omega_{\Q(i)} =
\Gamma(1/4)^2/(2\sqrt{2\pi})$ -- the \emph{lemniscate constant} -- is free of
$\sqrt{3}$, while $\Omega_{\Q(\omega)} = \Gamma(1/3)^3/(2\pi\sqrt{3})$ carries an
explicit $1/\sqrt{3}$ from the discriminant $D_{\Q(\omega)} = -3$.  This factor
propagates to odd-indexed critical values via the parity of the order-$6$ Hecke
character of $\Q(\omega)$.

\smallskip

This note is a companion to \cite{He25} (which studies $p$-adic stability of
Hecke $L$-values in anticyclotomic twist families) and to forthcoming work on the
Steinberg-edge obstruction for the $2$-adic $L$-function of $4.5.b.a$.


%% -------------------------------------------------------------------
\section{Preliminaries}
\label{sec:prelim}

\subsection{CM newforms and Gr\"{o}{\ss}encharacters}

Let $K = \Q(\sqrt{D_K})$ be an imaginary quadratic field with discriminant $D_K < 0$
and ring of integers $\OO_K$.  A \emph{Hecke Gr\"{o}{\ss}encharacter} $\psi$ of $K$
with infinity type $(k-1, 0)$ and conductor $\mathfrak{f} \subset \OO_K$ gives rise,
via the CM theta-series construction, to a newform
\[
  f = \theta(\psi) \in S_k\!\left(\Nm(\mathfrak{f})|D_K|,\, \chi_{D_K}\right)
\]
of weight $k$, level $\Nm(\mathfrak{f})|D_K|$, and nebentypus $\chi_{D_K}$.
We say $f$ has \emph{complex multiplication by $K$}, and write $\mathrm{CM}(f) = K$.

\subsection{The Chowla--Selberg formula}

The Chowla--Selberg formula \cite{CS67} expresses the CM period $\Omega_K \in \R_{>0}$
(the real period of the canonical CM elliptic curve over $K$, normalized by the Faltings
height) as a product of $\Gamma$-values:
\[
  \Omega_K \sim \prod_{a=1}^{|D_K|}
    \Gamma\!\left(\frac{a}{|D_K|}\right)^{\chi_{D_K}(a)/4h_K}
    \qquad (\text{up to an algebraic factor}),
\]
where $h_K = h(D_K)$ is the class number of $K$.
For the two class-number-$1$ fields of interest:
\begin{align}
  \label{eq:OmegaQi}
  \OmI &= \frac{\Gamma(1/4)^2}{2\sqrt{2\pi}}
    \qquad (K = \Q(i),\; D_K = -4), \\
  \label{eq:OmegaQomega}
  \OmW &= \frac{\Gamma(1/3)^3}{2\pi\sqrt{3}}
    \qquad (K = \Q(\omega),\; D_K = -3).
\end{align}
The period $\OmI$ is the \emph{lemniscate constant}; it satisfies
$\OmI^2 = \pi^{-1}\cdot\Gamma(1/4)^4/16$ and involves no $\sqrt{3}$.
The period $\OmW$ contains an explicit factor of $1/\sqrt{3}$, inherited
from $D_K = -3$.

\subsection{Damerell's theorem and Shimura algebraicity}

\begin{theorem}[Damerell \cite{Dam70, Dam71}; Shimura \cite{Shi76}]
\label{thm:damerell}
Let $f$ be a CM newform of weight $k$ with $\mathrm{CM}(f) = K$.  For each
critical integer $1 \leq m \leq k-1$, define $\alpha_m(f)$ by \eqref{eq:alpha}.
Then $\alpha_m(f) \in \overline{\Q}$.  More precisely, $\alpha_m(f)$ lies in
the abelian extension of $K$ cut out by the Gr\"{o}{\ss}encharacter $\psi$.
\end{theorem}

\begin{remark}
For weight-$5$ CM newforms with one-dimensional eigenspaces (as in the forms
studied here), Shimura's theorem \cite{Shi76} shows that the algebraicity field
is explicitly computable from the Hecke character data.  The Katz framework
\cite{Kat76, Kat78} provides a $p$-adic interpolation of these values via
the Maass--Shimura operator.
\end{remark}


%% -------------------------------------------------------------------
\section{Main theorem and conjecture}
\label{sec:main}

\subsection{The newform \texorpdfstring{$4.5.b.a$}{4.5.b.a}}

In LMFDB notation \cite{LMFDB}, the label $4.5.b.a$ denotes the unique CM newform of
weight $k=5$, level $N=4$, character $\chi_{-4}$ (the Kronecker symbol
$\left(\frac{-4}{\cdot}\right)$), with $\mathrm{CM}(\mathit{4.5.b.a}) = \Q(i)$.
Its key data (LMFDB-verified):
\begin{itemize}
  \item $a_2 = -4$ (Steinberg local type at $p=2$, since $N=4=2^2$);
  \item Hecke Gr\"{o}{\ss}encharacter $\psi_{\min}$ of conductor $(1+i)^2 \subset \Z[i]$,
        infinity type $(4,0)$;
  \item Self-minimal twist;
  \item Eta product: $f = \eta(z)^4\eta(2z)^2\eta(4z)^4$.
\end{itemize}

\begin{theorem}[Computational; PARI/GP, $80$-digit precision]
\label{thm:main}
Let $f = 4.5.b.a$ and $\OmI = \Gamma(1/4)^2/(2\sqrt{2\pi})$.  Then:
\begin{enumerate}
  \item[\rm(1)] \textup{[Damerell--Shimura, Theorem \ref{thm:damerell}]}
        The values $\alpha_m(f)$ are algebraic for $m = 1,2,3,4$.
  \item[\rm(2)] \textup{[Computational, PARI $80$-digit]}
        \begin{equation}\label{eq:ladder}
          (\alpha_1, \alpha_2, \alpha_3, \alpha_4) =
          \left(\frac{1}{10},\; \frac{1}{12},\; \frac{1}{24},\; \frac{1}{60}\right)
          \in \Q^4.
        \end{equation}
  \item[\rm(3)] \textup{[Ω-independent, consequence of \eqref{eq:ladder}]}
        All consecutive critical ratios are rational:
        \begin{equation}\label{eq:ratios}
          \frac{\pi \cdot L(f,1)}{L(f,2)} = \frac{6}{5},\quad
          \frac{\pi \cdot L(f,2)}{L(f,3)} = \frac{1}{2},\quad
          \frac{\pi \cdot L(f,3)}{L(f,4)} = \frac{2}{5}.
        \end{equation}
  \item[\rm(4)] \textup{[Comparative, PARI $80$-digit]}
        For the $\Q(\omega)$-CM forms $g = 27.5.b.a$ and $h = 12.5.c.a$:
        \begin{align*}
          (\alpha_1(g), \alpha_2(g), \alpha_3(g), \alpha_4(g))
            &= \left(\tfrac{27\sqrt{3}}{4},\; \tfrac{9}{4},\;
               \tfrac{\sqrt{3}}{4},\; \tfrac{1}{9}\right),\\
          (\alpha_1(h), \alpha_2(h), \alpha_3(h), \alpha_4(h))
            &= \left(2\sqrt{3},\; \tfrac{4}{3},\;
               \tfrac{2\sqrt{3}}{9},\; \tfrac{1}{9}\right).
        \end{align*}
        In both cases $\alpha_m \in \Q(\sqrt{3})\setminus\Q$ for $m$ odd,
        and $\alpha_m \in \Q$ for $m$ even.
        In particular, $R(g) = 3\sqrt{3}$ and $R(h) = 3\sqrt{3}/2$, both
        irrational.
\end{enumerate}
\end{theorem}

\begin{remark}[Parity dichotomy]
Table \ref{tab:ratios} summarizes the Ω-independent ratios.  The \emph{mixed-parity}
ratio $R(f) = \pi L(f,1)/L(f,2)$ is the sharp diagnostic: it is rational
exclusively for $K = \Q(i)$ among the three forms tested.  By contrast,
the \emph{even/even} ratio $\pi^2 L(f,2)/L(f,4)$ and \emph{odd/odd} ratio
$\pi^2 L(f,1)/L(f,3)$ are rational for \emph{all} three forms, because
the $\sqrt{3}$ factors cancel pairwise in the $\Q(\omega)$ cases.
\end{remark}

\begin{table}[h]
\centering
\begin{tabular}{lccc}
\hline
Ratio & $4.5.b.a$ & $27.5.b.a$ & $12.5.c.a$\\
\hline
$\pi\cdot L(f,1)/L(f,2)$ & $\mathbf{6/5}$ & $3\sqrt{3}$ & $3\sqrt{3}/2$\\
$\pi^2\cdot L(f,2)/L(f,4)$ & $5$ & $81/4$ & $12$\\
$\pi^2\cdot L(f,1)/L(f,3)$ & $12/5$ & $27$ & $9$\\
\hline
\end{tabular}
\caption{$\Omega$-independent ratios for three CM weight-$5$ newforms.
Mixed-parity ratio (top row) is $\in\Q$ only for $K=\Q(i)$.}
\label{tab:ratios}
\end{table}

\begin{conjecture}[Scope extension]
\label{conj:scope}
Let $K = \Q(\sqrt{-d})$ be imaginary quadratic with $h_K = 1$ and let $f$ be a
CM weight-$5$ newform with $\mathrm{CM}(f) = K$.
\begin{enumerate}
  \item[\rm(a)] $R(f) = \pi L(f,1)/L(f,2) \in \Q \iff K = \Q(i)$.
  \item[\rm(b)] $K = \Q(\omega) \Rightarrow R(f) \in \Q(\sqrt{3})\setminus\Q$.
  \item[\rm(c)] \textup{[Predicted]} $K = \Q(\sqrt{-d})$ for $d \neq 1,3$ implies
        $R(f) \in \Q(\sqrt{d})\setminus\Q$.
\end{enumerate}
Parts \textup{(a)} and \textup{(b)} are supported by three PARI $80$-digit
data points (Theorem \ref{thm:main}).
Part \textup{(c)} is a structural prediction from the Chowla--Selberg heuristic
(\S\ref{sec:structure}); computational verification for $d \in \{7, 11\}$
is ongoing.
\end{conjecture}


%% -------------------------------------------------------------------
\section{Computational verification}
\label{sec:comp}

\subsection{Setup}

We use PARI/GP 2.15.4 \cite{PARI}.  For each form, we:
\begin{enumerate}
  \item Call \texttt{mfinit([N, k, chi], 1)} (new subspace, character $= \chi_D$);
  \item Extract \texttt{mfeigenbasis(mf)} and identify the unique eigenform
        matching LMFDB Fourier coefficients $a_n$ for $n = 2, \ldots, 8$
        (disambiguation against old-form lifts);
  \item Compute $L(f, m)$ via \texttt{lfunmf} at $80$-digit precision;
  \item Normalize by $\Omega_K^{k-1}/\pi^{k-1-m}$, computed from
        PARI's \texttt{gamma(1/4)} or \texttt{gamma(1/3)} at $80$ digits;
  \item Reconstruct rationals via \texttt{bestappr(x, 10\textasciicircum{}20)}.
\end{enumerate}

\subsection{Verification table for \texorpdfstring{$4.5.b.a$}{4.5.b.a}}

\begin{table}[h]
\centering
\begin{tabular}{clll}
\hline
$m$ & $L(f,m)$ (leading digits) & $\alpha_m$ (closed form) & Residual\\
\hline
$1$ & $0.387393\ldots$ & $1/10$ & $< 3\times10^{-76}$\\
$2$ & $0.075171\ldots$ & $1/12$ & $< 1\times10^{-77}$\\
$3$ & $0.005943\ldots$ & $1/24$ & $< 5\times10^{-76}$\\
$4$ & $0.000370\ldots$ & $1/60$ & $< 2\times10^{-60}$\\
\hline
\end{tabular}
\caption{Damerell ladder for $f = 4.5.b.a$; $\OmI = \Gamma(1/4)^2/(2\sqrt{2\pi})$.
Residuals are $|\alpha_m - \text{closed form}|$ from PARI computation.}
\label{tab:ladder}
\end{table}

\begin{remark}
The slightly larger residual at $m=4$ reflects the smaller absolute value of
$L(f,4)$ (fewer contributing terms in PARI's series acceleration).  All
residuals are well within $80$-digit precision; true residuals are zero.
\end{remark}

\subsection{Disambiguation protocol}

An earlier sweep (without disambiguation) produced false positives for higher-level
forms, because PARI's \texttt{mfinit} at level $N > 4$ with character $\chi_{-4}$
returns bases that include old-form lifts of $4.5.b.a$.  The LMFDB trace-matching
step (checking $a_n$ for $n = 2,\ldots,8$) is essential to identify the correct
newform within the eigenspace basis.


%% -------------------------------------------------------------------
\section{Structural explanation via Chowla--Selberg}
\label{sec:structure}

\subsection{Period arithmetic}

By \eqref{eq:OmegaQi} and \eqref{eq:OmegaQomega}:
\begin{align}
  \OmI^4 &= \frac{\Gamma(1/4)^8}{16\cdot(2\pi)^2}
           = \frac{\Gamma(1/4)^8}{64\pi^2} \in \Q\cdot \Gamma(1/4)^8/\pi^2, \\
  \OmW^4 &= \frac{\Gamma(1/3)^{12}}{16\pi^4\cdot 9}
           = \frac{\Gamma(1/3)^{12}}{144\pi^4} \in \Q\cdot \Gamma(1/3)^{12}/\pi^4.
\end{align}
The key contrast is:
\begin{itemize}
  \item $\OmI^4$ is a pure product of $\Gamma(1/4)^8$ and rational powers of $\pi$.
        No $\sqrt{3}$ appears anywhere.
  \item $\OmW^4$ also contains no explicit $\sqrt{3}$ -- \emph{but} the four-th power
        of $1/\sqrt{3}$ in $\OmW$ equals $1/9 \in \Q$, so the $\sqrt{3}$ is invisible
        at even powers of $\OmW$.  At odd multiples (such as in $\OmW^1, \OmW^3, \ldots$),
        the factor $1/\sqrt{3}$ remains irrational.
\end{itemize}

More precisely, $\alpha_m = L(f,m)/(\pi^{k-1-m}\OmK^{k-1})$ involves $\OmK^{k-1}
= \OmK^4$.  For K = Q(ω), if one writes $L(f,m) = c_m \cdot \OmW^{k-1} \cdot \pi^{k-1-m}$
with $c_m \in \overline{\Q}$, then the question is whether $c_m \in \Q$.  The Hecke
character parity of $\Q(\omega)$ places a residual $\sqrt{3}$ at odd $m$ (see below).

\subsection{Hecke character parity for \texorpdfstring{$\Q(\omega)$}{Q(omega)}}

For $K = \Q(\omega)$, the unit group $U_K = \OO_K^\times = \{\pm 1, \pm\omega, \pm\omega^2\}$
has order $6$.  The Hecke Gr\"{o}{\ss}encharacter $\psi$ for a weight-$k$ CM form is
of the form $\psi(\mathfrak{a}) = \alpha^{k-1}$ for a generator $\alpha \in K^\times$
of $\mathfrak{a}$.  For odd $m$, the period decomposition of $L(\psi|\cdot|^{1-m}, 0)$
involves the imaginary part of the fundamental period $\omega = e^{2\pi i/3}$, namely
$\mathrm{Im}(\omega) = \sqrt{3}/2$.  At odd $m$, one factor of $\sqrt{3}/2$ remains
unpaired:
\[
  \left(\frac{\sqrt{3}}{2}\right)^{2j} = \frac{3^j}{4^j} \in \Q,
  \qquad
  \left(\frac{\sqrt{3}}{2}\right)^{2j+1} = \frac{3^j\sqrt{3}}{4^{j+1}} \notin \Q.
\]
This is the mechanism by which $\alpha_m \in \Q(\sqrt{3})\setminus\Q$ for odd $m$
when $K = \Q(\omega)$.

\subsection{Absence of \texorpdfstring{$\sqrt{3}$}{sqrt(3)} contamination for \texorpdfstring{$K=\Q(i)$}{K=Q(i)}}

For $K = \Q(i)$, the unit group $U_{\Q(i)} = \{\pm 1, \pm i\}$ has order $4$.
The fundamental period is $\OmI = \Gamma(1/4)^2/(2\sqrt{2\pi})$, which is real
and contains no $\sqrt{3}$.  The Hecke character $\chi_{-4}$ is the Kronecker symbol
for $D = -4$: it takes values in $\{\pm 1\}$ only.  There is no analogue of the
$\mathrm{Im}(\omega) = \sqrt{3}/2$ pairing; at all critical $m$, the period
decomposition involves only rational combinations of $\OmI$ and powers of $\pi$,
with coefficients in $\Q$.  Hence all four $\alpha_m \in \Q$.

\subsection{Heuristic for general \texorpdfstring{$K$}{K}}

For $K = \Q(\sqrt{-d})$ with $h_K = 1$ and $d \notin \{1, 3\}$, the Chowla--Selberg
formula \eqref{eq:OmegaQi} introduces $\Gamma$-values at fractions $a/d$ with
$1 \leq a \leq d$, $\gcd(a,d) = 1$.  The resulting period $\Omega_K$ carries
algebraic factors from $\Q(\sqrt{-d})$ (via the character sum $\sum \chi_{D_K}(a)\log
\Gamma(a/d)$).  At odd critical indices $m$, these algebraic factors do not cancel
and yield $\alpha_m \in \Q(\sqrt{d})\setminus\Q$, consistent with Conjecture
\ref{conj:scope}(c).

\begin{remark}[\textbf{Open: formal proof for $\Q(\omega)$}]
\label{rem:TBD}
A complete proof that $\alpha_1, \alpha_3 \in \Q(\sqrt{3})\setminus\Q$ for
$\Q(\omega)$-CM forms would follow from:
\begin{enumerate}
  \item[\rm(i)] Shimura \cite{Shi76} \S{}5, giving explicit period decompositions
        for CM Hecke $L$-functions;
  \item[\rm(ii)] Damerell \cite{Dam71} \S{}4, giving explicit denominators of $\alpha_m$
        at conductor $(1+i)^2$;
  \item[\rm(iii)] Elementary $\sqrt{3}$-arithmetic: $(\sqrt{3}/2)^{2j+1} \notin \Q$ for $j \geq 0$.
\end{enumerate}
Formalizing step (i)--(ii) in the weight-$5$ setting is left as a \emph{[TBD: prove]}
marker; the present note establishes the result computationally.
\end{remark}


%% -------------------------------------------------------------------
\section{Conclusion}

We have shown computationally that the CM newform $f = 4.5.b.a$ (weight $5$, level $4$,
$\mathrm{CM}(f) = \Q(i)$) has a fully rational Damerell ladder, $\alpha_m \in \Q$ for
all $m = 1,2,3,4$, with the clean closed-form values \eqref{eq:ladder}.  The ratio
$R(f) = \pi L(f,1)/L(f,2) = 6/5$ is an Ω-independent rational invariant unique to
$K = \Q(i)$ among the forms tested.

The structural reason is the lemniscate character of $\Omega_{\Q(i)}$: free of $\sqrt{3}$,
purely a Gamma-value product at $1/4$.  The $\Q(\omega)$-CM forms exhibit a
rationality dichotomy between even and odd critical indices, with $\sqrt{3}$
contamination at odd $m$ traceable to $\mathrm{Im}(\omega) = \sqrt{3}/2$.

The invariant $R(f) = 6/5$ is, to our knowledge, new.  It provides a sharp
Ω-independent criterion distinguishing $K = \Q(i)$ from $K = \Q(\omega)$ at the
level of special $L$-values, complementing the Steinberg-edge $v_2$-fingerprint
studied in related work.


%% -------------------------------------------------------------------
\section*{Acknowledgements}

Computations performed with PARI/GP 2.15.4.  LMFDB data accessed at
\texttt{www.lmfdb.org}.


%% -------------------------------------------------------------------
\begin{thebibliography}{99}

\bibitem{CS67}
S.~Chowla and A.~Selberg,
\textit{On Epstein's zeta-function},
J. Reine Angew. Math.\ \textbf{227} (1967), 86--110.
%% [TBD: confirm pages; standard reference, classical]

\bibitem{Dam70}
R.~M.~Damerell,
\textit{$L$-functions of elliptic curves with complex multiplication, {I}},
Acta Arith.\ \textbf{17} (1970), 287--301.
%% [TBD: live-verify pages via zbMATH]

\bibitem{Dam71}
R.~M.~Damerell,
\textit{$L$-functions of elliptic curves with complex multiplication, {II}},
Acta Arith.\ \textbf{19} (1971), 311--317.
%% [TBD: live-verify pages via zbMATH]

\bibitem{Shi76}
G.~Shimura,
\textit{The special values of the zeta functions associated with cusp forms},
Comm.\ Pure Appl.\ Math.\ \textbf{29} (1976), 783--804.
%% [TBD: live-verify via CrossRef]

\bibitem{Kat76}
N.~M.~Katz,
\textit{$p$-adic interpolation of real analytic Eisenstein series},
Ann.\ Math.\ \textbf{104} (1976), 459--571.
%% [TBD: live-verify via CrossRef]

\bibitem{Kat78}
N.~M.~Katz,
\textit{$p$-adic $L$-functions for {CM} fields},
Invent.\ Math.\ \textbf{49} (1978), 199--297.
DOI: \href{https://doi.org/10.1007/BF01390187}{10.1007/BF01390187}.
%% CONFIRMED: DOI live-verified in M57 session (2026-05-06)

\bibitem{CO77}
H.~Cohen and J.~Oesterlé,
\textit{Dimensions des espaces de formes modulaires},
in: \textit{Modular Functions of One Variable VI},
Lecture Notes in Math.\ \textbf{627}, Springer, Berlin, 1977.
%% [TBD: confirm exact pages for dimension formula]

\bibitem{LMFDB}
The LMFDB Collaboration,
\textit{The $L$-functions and Modular Forms Database},
\url{http://www.lmfdb.org} (2024).
%% Entry 4.5.b.a confirmed 2026-05-06: weight 5, level 4, a_2=-4, CM Q(i), self-minimal.

\bibitem{He25}
W.~He,
\textit{Stability of $p$-adic valuations of {H}ecke $L$-values in anticyclotomic
twist families},
Math.\ Ann.\ \textbf{392} (2025), 399--468.
arXiv:\href{https://arxiv.org/abs/2308.15051}{2308.15051}.
%% CONFIRMED: arXiv 2308.15051, Math. Ann. confirmed in M55 session + this session.

\bibitem{PARI}
The PARI~Group,
\textit{{PARI/GP} version 2.15.4},
Bordeaux, 2023,
\url{http://pari.math.u-bordeaux.fr/}.

\end{thebibliography}

\end{document}
